\section{Discussion}\label{s:discussion}

The signature check worked when its reference held every declared workload. With no
example of hidden training in its fit, $s_R$ flagged every hidden-training group, and it
raised no alarm on test-day decode (\cref{s:works}). When the reference missed a declared
workload, the check failed in one of two ways (\cref{s:monitoring}). A workload missing
from both parts of the reference kept the monitor in alarm almost continuously while it
ran. This failure was loud, so the verifier could see that something was wrong. A workload in
the threshold set but not in the fit raised $\tau_\alpha$ and blinded the monitor. That
monitor raised no more alarms on declared work than the working one, so the failure was
silent. Such a workload could also mask a loud failure. Training in the reference also blinded the
monitor. So a quiet monitor is no evidence of a working one. The signature check needs no examples of hidden
training in its fit, but the verifier still needs them as positive controls, to test that
the monitor works.

\paragraph{Acquiring a reference and auditing the monitor.} The signature check is only as good as
its reference. The reference is recorded on the monitored chip, where the operator
controls the work. The verifier must therefore control how it is recorded.
\begin{itemize}
  \item \emph{Witness it.} Our reference recordings are our own honest runs. A real
        verifier needs recordings during which the operator cannot train. Both parts of
        the reference are targets. Training at a share of about $\alpha$ in the threshold
        set takes detection of the unthrottled run to zero, and full fine-tuning at
        \EfourShareHonest\% of the fit hides full fine-tuning (\cref{s:monitoring}).
  \item \emph{Span sessions.} A reference from one session flags \SessOther\% of decode
        from the other days (\cref{s:declaration}).
  \item \emph{Test with held-out declared workloads, not only held-out attacks.} A
        declared workload missing from the reference raises alarms
        (\cref{s:declaration}), unless a workload in the threshold set but not in the
        fit masks it (\cref{s:monitoring}).
  \item \emph{Audit the monitor with positive controls that sit close to the
        reference.} The silent monitor flags \TfourPosCtrl\% of honest training, so one
        control run exposes it. But a monitor that flags a control can still miss other hidden work. At the share
        of training in the threshold set that hides the unthrottled run, honest training
        is still flagged (\cref{s:monitoring}). Honest training sits far from the
        reference, so it cannot expose a raised threshold. A control must sit as close
        to the reference as the hidden work it represents.
\end{itemize}
Anomaly detection is hard when normal behaviour is open-ended \citep{sommer2010}. An
agreement can close it, by requiring the operator to declare every workload, so that
anything else is an anomaly by rule. The reference then plays the role of the template in warhead verification, where
the hard problem is to show that the template item is authentic \citep{yan2015}. A
classifier fitted elsewhere, such as the hardened forest of \citet{gargiulo2026}, might
screen recordings offered as declared work, but that forest also flags \GkHardPrefill\%
of our declared prefill (\cref{s:missed}).

\paragraph{What an alarm means.} An alarm from $s_R$ says only that the GPU ran work
unlike its declaration. $s_R$ asks ``is this declared work?'', not ``is this
training?''. Its score does not rank the work it flags, because co-location, honest
training and non-AI work share one band (\cref{f:f7}). As a monitor it is in alarm for a fraction
\ARoneBusy{} of the time on undeclared non-AI work, where the classifier, $s_C^{\mathrm{dev}}$, stays silent
(\cref{t:t2}). Nor does an alarm show that a model was trained.
% results/ceiling/s2.json: tests.r1.primary, rate_a = rate_b = 1.0
The signature check flags every window of a run that discards every update, as it does
for training that applies its updates. An alarm therefore works like an anomaly under
nuclear safeguards~\citep[\S\S 11.24--11.26]{iaea2022glossary}. It opens an inquiry, in
which the operator explains the departure or the verifier inspects. No conclusion is drawn until the
operator has had the chance to explain~\citep[Art.~4.d]{iaea1997ap}. The false-alarm rate
per hour sets how often the verifier must inquire, and so the cost of the regime
(\cref{s:declaration}).

\paragraph{Limits.}
\begin{itemize}
  \item \emph{Scale.} We use one model and one GPU, with one run per co-location share,
        and we do not test one probe on several GPUs at once. Co-location is development
        data from one session. The evasions did not adapt to the detector. Our corpus
        has no diluted LoRA. The scale, the chip activity measured and the threat model
        of \citet{scher2025} are wider than ours.
  \item \emph{Code access.} \citet{scher2025} rely on temporary code access to
        label the recordings, and to review the code after an alarm. We test neither.
        Our reference recordings are our own honest runs, and we test only whether the
        signal departs from them.
  \item \emph{Hardware.} We test no transfer between chips. Every detector we fit,
        the classifiers included, is fitted and scored on the GPU we monitor, and the
        published forests of \citet{gargiulo2026} come from the same GPU model and probe
        model. The default approach must transfer in practice, so our comparison
        favours it. The signature check needs no transfer, but its reference must hold
        across sessions.
  \item \emph{Exposure.} The test day holds no prefill and only \ExpDecodeH~h of
        declared decode, so the false-alarm bounds are loose, and declared prefill is a
        risk we could not measure on held-out data.
  \item \emph{Amplitude.} The unthrottled run draws a median AC RMS current of
        \AcRmsUnpaced~A, against \AcRmsDecode~A for development decode, and the per-run interquartile ranges do not overlap. So
        amplitude alone would separate this run from decode, and the run is not a hard
        case for a detector that uses amplitude. Our features are amplitude-free and do
        not use it. A harder test would match the amplitude of declared decode.
  \item \emph{Design.} Our results depend on our split of the reference into a fit and
        a threshold set (\cref{t:t7}), whereas \citet{scher2025} specify one baseline.
        One short calibration run each of decode and idle sets $\tau_\alpha$
        (\cref{t:corpus}). That idle run is unlike test idle, since $s_R$ flags
        \RoneCalIdle\% of it but only \FprRoneIdle\% of test idle. Resampling whole calibration runs moves
        $\tau_\alpha$ over [\TauRoneLo, \TauRoneHi].
  \item \emph{Training data in view.} The signature check uses no training or evasion
        examples in its reference, but three earlier choices drew on them. Its features
        come from \citet{gargiulo2026}, who chose them to separate training from
        inference. When we chose its design, we required each candidate to flag at least
        \PosControlMin\% of honest-training windows from an unseen session. Every
        candidate passed, so this requirement did not decide the choice. We also measured its detection on development data
        before we scored any test data. Only its test false-positive rates and the monitor
        results are held out.
  \item \emph{Contamination.} We test contamination only in the exploratory checks of
        \cref{s:monitoring}, with training drawn from our own runs.
  \item \emph{Measurement chain.} Our measurement chain and that of \citet{gargiulo2026}
        differ in sample rate, filtering and feature set (\cref{a:gk}).
\end{itemize}
